\begin{abstract}
Recovering fluorescent structures that overlap in projection requires measurements that distinguish their axial arrangements. We learn depth-encoding illumination for fixed-focus fluorescence microscopy using the optical model alone, without specimen examples or a reconstruction decoder. A fixed, defocused phase plate produces depth-dependent excitation as the illumination angle is scanned, while the specimen and detection focus remain stationary. We optimize its phase to minimize squared correlations between camera responses from different depths at the same lateral position. In simulations of mouse pyramidal neurons, the learned plate improves axial recovery with total variation, two plug-and-play priors and a feed-forward decoder at eightfold measurement compression and matched nominal excitation dose. With total variation, mean axial-profile Wasserstein-1 error is $0.22\um$, compared with $0.37\um$ for the untrained plate. Plates optimized using depth Fisher information or worst-case coherence achieve $0.25$ and $0.26\um$, respectively, under the same optimization budget. Mean iterative errors are close to those of an end-to-end plate, whose more extensively trained decoder performs better. Controlled crossings also show improved filament separation over the untrained plate with total variation and pretrained plug-and-play reconstruction. Decoder-free optical learning thus allows illumination to be optimized before a reconstruction method is chosen.
\end{abstract}
